\scanpage{42}
\mspage{- 15 -}
\margin{n°357}

déterminée par cette même norme ; son dual est alors $\mathcal M_b^1(E)$, espace des mesures bornées sur $E$.

PROPOSITION 5.- \uline{Soit} $G$ \uline{un groupe localement compact opérant sur un espace} $E$ \uline{localement compact} ; \uline{pour tout} $x\in G$, \uline{soit} $\varphi_x$ \uline{l'automorphisme} $z\to xz$ \uline{déterminé par} $x$ \uline{dans} $E$. \uline{Alors les applications} $(x,f)\to f\circ\varphi_x$ \uline{de} $G\times\mathcal K(E)$ \uline{dans} $\mathcal K(E)$, \uline{et de} $G\times\overline{\mathcal K(E)}$ \uline{dans} $\overline{\mathcal K(E)}$, \uline{sont continues}.

Soit $V$ un voisinage de $0$ dans $\mathcal K(E)$ ; il y aura donc, pour tout compact $K$ dans $E$, un $\delta(K)>0$ tel que $f\in V$ chaque fois que $\sup_E|f|\le\delta(K)$ et que le support de $f$ est contenu dans $K$. Soient $x_0$ un point de $G$ et $K_0$ un voisinage compact de $x_0$ dans $G$. Soit $W$ l'ensemble des $g\in\mathcal K(E)$ tels que la fonction $z\to 2g(xz)$ appartienne à $V$ quel que soit $x\in K_0$ ; c'est un voisinage de $0$ dans $\mathcal K(E)$, car $W$ est convexe et symétrique, et, si $K$ est compact, tout $g$ de support contenu dans $K$ et tel que $\sup_E|g|\le\delta(K_0^{-1}K)/2$ est dans $W$. Soit $f_0\in\mathcal K(E)$ ; soit $K'$ son support ; soit $K''=K_0^{-1}K'$ ; pour tout $x\in K_0$, le support de $z\to f_0(xz)$ est contenu dans $K''$ ; la fonction $(x,z)\to f_0(xz)$ étant uniformément continue sur $K_0\times K''$, il y a un voisinage $U$ de $x_0$ dans $K_0$ tel que l'on ait
\[
|f_0(xz)-f_0(x_0z)|\le\delta(K'')/2
\]
pour tout $x\in U$ et tout $z\in K''$ ; donc, pour $x\in U$, la fonction
\[
z\to 2f_0(xz)-2f_0(x_0z)
\]
appartient à $V$. Comme $V$ est convexe, on en conclut que, si l'on pose $f=f_0+g$, la fonction $f\circ\varphi_x-f_0\circ\varphi_{x_0}$ est dans $V$ quels que soient $x\in U$ et $g\in W$, ce qui démontre la première partie de la proposition. La démonstration de la seconde est analogue mais plus simple, et est laissée au lecteur pour son amusement.

COROLLAIRE.- \uline{Soit} $G$ \uline{un groupe localement compact opérant sur un espace loca-}

\margin{52}
